\section{Periodic concentration estimates}\label{sec:periodic}

The boundary argument requires quantitative information about the
spatial spectrum, including the change of its leading eigenvalue when
the inverse temperature doubles.  We extract that information from
finite periodic-chain estimates in~\cite{periodic}.  The extraction is
given in detail because the periodic gap theorem alone does not provide
these boundary estimates.  Throughout, terminating decimals denote
exact rational numbers, and eigenvalue lists count multiplicities.

For even $n\ge4$, put
\begin{equation}\label{per:purities}
 S(n,b)=\frac{Z_{2n}(b)}{Z_n(b)^2},\qquad
 T(n,b)=\frac{Z_n(2b)}{Z_n(b)^2}.
\end{equation}
The first is the purity of the spatial probabilities
$|\mu_i(b)|^n/Z_n(b)$, where $(\mu_i(b))$ is the nonzero spectrum of
$X_b$; the second is the purity of the physical Gibbs probabilities.
Both lie in $(0,1]$.  Doubling $n$ squares the first probability list,
whereas doubling $b$ squares the second.

\begin{lemma}[Squaring and concentration]\label{per:purity}
Let $(w_i)$ be a finite or countable probability list, and suppose its
purity is at least $1-u$, where $0\le u<1$.  Normalizing the squares
$w_i^2$ gives a list with purity defect at most
\begin{equation}\label{per:f}
 f(u)=\frac{u^2}{2(1-u)^2}.
\end{equation}
If $u<1/2$, the largest weight is unique and at least $1-u$.  Writing
$x$ for the sum of all the other weights, one has
\begin{equation}\label{per:concentration-elementary}
 x\le\frac{1-\sqrt{1-2u}}2.
\end{equation}
In particular, for $0\le u\le .1$,
\begin{equation}\label{per:d}
 \frac{x}{1-x}\le d(u),\qquad d(u)=\frac{u}{2-3u}.
\end{equation}
\end{lemma}

\begin{proof}
Set $P=\sum_iw_i^2$ and let $P'$ be the purity after squaring.
Nonnegativity justifies rearrangement of the countable sums, and
\[
 1-P'=\frac{2\sum_{i<k}w_i^2w_k^2}{P^2}
 \le\frac{2(\sum_{i<k}w_iw_k)^2}{P^2}
 =\frac{(1-P)^2}{2P^2}\le f(u).
\]
Here $f$ is increasing on $[0,1)$.  A summable nonzero nonnegative list
attains its maximum, and $P\le\max_iw_i$.  If $u<1/2$, the maximum
therefore exceeds $1/2$ and is unique.  Thus $x\le u<1/2$, and
\[
 2x(1-x)\le 1-P\le u,
\]
which gives~\eqref{per:concentration-elementary}.  Also
$x\le u/[2(1-x)]\le u/[2(1-u)]$; applying the increasing function
$x\mapsto x/(1-x)$ proves~\eqref{per:d}.
\end{proof}

For later reference we also record explicitly the interaction of the
two periodic purities.  This is \cite[Lemma 4.2]{periodic}; we include its short proof by
partition-function cancellation.

\begin{lemma}[Coupled update]\label{per:coupled}
Suppose that $n\ge4$ is even, $b>0$, and
$S(n,b)\ge s$, $T(2n,b)\ge t$, with $0<s,t\le1$.
Set
\[
 p_+=f(1-s^2t).
\]
If $p_+<1$, then
\begin{equation}\label{per:coupled-formula}
 1-S(2n,2b)\le p_+,\qquad
 1-T(4n,2b)\le f\bigl(1-t^2(1-p_+)\bigr).
\end{equation}
Either upper defect bound may be increased before subsequent use,
provided it stays below one.
\end{lemma}

\begin{proof}
Directly from~\eqref{per:purities},
\[
 S(n,2b)=S(n,b)^2\frac{T(2n,b)}{T(n,b)^2}
 \ge s^2t.
\]
Apply spatial squaring and Lemma~\ref{per:purity}.  A second cancellation
then gives
\[
 T(4n,b)=T(2n,b)^2\frac{S(2n,2b)}{S(2n,b)^2}
 \ge t^2(1-p_+).
\]
Physical squaring proves the second inequality.  Monotonicity of $f$
justifies increasing either defect bound.
\end{proof}

We next state the finite inputs gathered in the companion
\cite[Corollary 5.6]{periodic}. Its shift, spin normalization and unit
bilinear coupling are exactly those used here.
In particular, we use an intermediate row in the initialization proof,
not merely its final conclusion.

\begin{proposition}[Finite periodic inputs from~\cite{periodic}]
\label{per:native}
Set
\begin{equation}\label{per:constants}
 \ell=.999,\quad U=1.00139,\quad V=.872,\quad
 m_J=1.0657205,\quad m_N=.2783155.
\end{equation}
At $b_0=49/4$, the spatial eigenvalue list decomposes as one list $J$
and three identical lists $N$. Here $J$ is the spectrum on the subspace
fixed by every $U_g$, $g\in D_2$, and $N$ is the common spectrum on
each of the other three simultaneous eigenspaces of these commuting
operators.
In $J$ the absolute eigenvalues are at
most $U$, and their twelfth powers sum to at most $m_J$; in each $N$
they are at most $V$, and their twelfth powers sum to at most $m_N$.
There is an eigenvalue in $J$ whose modulus exceeds $\ell$.
Furthermore,
\begin{equation}\label{per:native-first-row}
 1-S(72,b_1)\le\frac{151249}{2500000}=.0604996.
\end{equation}
For the periodic chains of $72$ and $120$ sites, the shifted ground
energy is strictly negative.  Consequently, at each of these lengths
and every $b>0$, an individual unnormalized ground Boltzmann weight is
greater than one.
\end{proposition}

No boundary Hamiltonian occurs in Proposition~\ref{per:native}.
The finite periodic thermal and variational certificates supporting it
are given in the cited source.  We now deduce all the spatial estimates
that will be used for boundaries. The residual ratios will contract
the boundary-weighted tails. The upper and lower bounds comparing
successive leading eigenvalues will control the change of temperature
in the boundary concentration and physical-purity estimates,
respectively. Finally, positivity and rotation invariance of the
leading line will control the same-field pairing at odd lengths.

\begin{proposition}[Concentration at the doubling scales]
\label{per:concentration}
For every $j\ge0$, $X_{b_j}$ has a unique simple eigenvalue of largest
modulus.  This eigenvalue is positive; denote it by
$\lambda_j=\|X_{b_j}\|$.  Its eigenspace is fixed pointwise by
$\mathcal G$.  Let $\rho_j$ be the supremum of the absolute values of
all other eigenvalues divided by $\lambda_j$.  Then
\begin{equation}\label{per:initial-ratios}
 \ell<\lambda_0\le U,\qquad
 \rho_0<.874,\qquad \rho_1<.9525,\qquad \rho_2<.971.
\end{equation}
For $j\ge2$, define
\begin{equation}\label{per:scales}
 n_j=120\,2^{j-2},\qquad M_j=2n_j,\qquad
 u_2=.052,\qquad u_{j+1}=6u_j^2.
\end{equation}
The following bounds hold for every $j\ge2$:
\begin{align}
 1-S(n_j,b_j)&\le u_j,&
 1-T(M_j,b_j)&\le u_j, \label{per:tracked-purities}\\
 \rho_j^{n_j}&\le d(u_j). &&\label{per:spectral-ratio}
\end{align}
Finally, put $r_j=\lambda_{j+1}/\lambda_j^2$.  Then
\begin{equation}\label{per:lambda-ratios}
 0<r_j\le1\quad(j\ge0),\qquad
 r_j^{M_j}\ge(1-u_j)(1-u_{j+1})\quad(j\ge2).
\end{equation}
\end{proposition}

\begin{proof}
Until positivity has been proved, write $\lambda_j=\|X_{b_j}\|$ for
the maximal spectral modulus.

\emph{The initial spatial spectrum.}
Remove an eigenvalue in $J$ whose modulus exceeds $\ell$.
The remaining twelfth-power mass in $J$ is at most
$m_J-\ell^{12}<\ell^{12}$, while $V<\ell$ bounds each of the other
three sectors.  The removed eigenvalue is therefore the unique simple
eigenvalue of maximal modulus.  The exact rational inequalities
\[
 m_J-\ell^{12}<(.874\ell)^{12},\qquad V<.874\ell
\]
prove $\rho_0<.874$.

We need the total residual moments as well as this individual ratio
bound.  For every real $k\ge12$, define
\begin{equation}\label{per:D}
 D(k)=\frac{(m_J-\ell^{12})^{k/12}
       +3\bigl[V^k+(m_N-V^{12})^{k/12}\bigr]}{\ell^k}.
\end{equation}
The sum of the $k$th powers of all nondominant absolute eigenvalue
ratios at $b_0$ is at most $D(k)$.  Indeed,
$V^{12}<m_N<2V^{12}$; for a nonnegative list with entries at most
$V^{12}$ and total mass at most $m_N$, convexity of $x^{k/12}$ bounds
the sum by $(V^{12})^{k/12}+(m_N-V^{12})^{k/12}$.  This follows by
moving mass to fill one entry before the next; finite sublists followed
by a limit give the countable case.  Without an entry cap, the
corresponding bound is simply the total mass to the power $k/12$.
Applying these facts to the four sectors proves the claim.  In
particular, for even $k\ge12$,
\begin{equation}\label{per:Z-initial}
 Z_k(b_0)\le U^k(1+D(k)).
\end{equation}

\emph{Reseeding at $(120,b_2)$.}
Our next goal is to make both $1-S(120,b_2)$ and $1-T(240,b_2)$
at most $.052$, so that one quadratic recurrence will control them.
A coupled update will give these bounds from spatial and physical
seeds at $(60,b_1)$ and $(120,b_1)$, respectively. To obtain the
spatial seed, we combine a leading-eigenvalue lower bound from length
$72$ with a residual-moment upper bound at length $32$, then increase
the moment order to $60$. The same residual estimate also supplies
the ratio $\rho_1$ needed in the first boundary updates.

By~\eqref{per:native-first-row} and Lemma~\ref{per:purity}, the spatial
list at $(72,b_1)$ has a unique largest weight.  Thus $X_{b_1}$ also
has a unique simple eigenvalue of maximal modulus.  Moreover,
\[
 .0604996<2(.032)(.968)
\]
implies that this weight exceeds $.968$.  The periodic trial at length
$72$ gives $Z_{72}(b_1)>1$, so
\begin{equation}\label{per:lambda1}
 \lambda_1^{72}>.968.
\end{equation}
Physical squaring at length $32$ gives
$Z_{32}(b_1)\le Z_{32}(b_0)^2$.  Subtracting the dominant spatial
contribution and using~\eqref{per:Z-initial} and~\eqref{per:lambda1}
therefore yields
\begin{equation}\label{per:tail32}
 \sum_{i\ne0}\left|\frac{\mu_i(b_1)}{\lambda_1}\right|^{32}
 \le U^{64}(1+D(32))^2(.968)^{-32/72}-1
 <.9525^{32}.
\end{equation}
This proves $\rho_1<.9525$.  For any nonnegative list $(a_i)$ and
$q\ge p>0$, one has
$\sum_i a_i^q\le(\sum_i a_i^p)^{q/p}$ whenever the right side is finite.
Applying this to~\eqref{per:tail32} gives a residual $60$th moment below
$.9525^{60}$.  Retaining only the dominant term in the numerator of
$S(60,b_1)$ proves
\begin{equation}\label{per:seed-s}
 S(60,b_1)\ge s:=(1+.9525^{60})^{-2}.
\end{equation}

For the physical spectrum at length $120$, choose one ground
Boltzmann weight at $b_0$ and divide every other weight by it.
The chosen weight is greater than one by Proposition~\ref{per:native}.
Hence the sum of the other ratios is bounded above by
\[
 R=U^{120}(1+D(120))-1.
\]
When $b$ doubles, these ratios square, so their sum is at most $R^2$.
Keeping the chosen ground weight in the purity numerator now gives
\begin{equation}\label{per:seed-t}
 T(120,b_1)\ge t:=(1+R^2)^{-2}.
\end{equation}
This reasoning neither assumes ground-state simplicity nor discards
multiplicities.

Here are rational checks sufficient for every numerical comparison
just used and for the ensuing seed:
\begin{equation}\label{per:scalar-checks}
\begin{array}{rcl}
 D(32)&<&.0442,\\
 D(120)&<&.00000025,\\
 (.9856)^9&<&(.968)^4,\\
 U^{64}(1+.0442)^2/.9856-1&<&.2092<.9525^{32},\\
 s&>&.9002,\qquad t>.9372,\\
 f(1-.9002^2\,.9372)&<&.0502,\\
 f(1-.9372^2(1-.0502))&<&.0198.
\end{array}
\end{equation}
All comparisons are between explicit rational numbers after bounding
the cube roots in $D(32)$ by
\[
 (m_J-\ell^{12})^{1/3}<.426635,\qquad
 (m_N-V^{12})^{1/3}<.439736.
\]
Cubing verifies these two bounds; raising the bounds to the eighth
power gives the displayed bound on $D(32)$.  For $D(120)$ no roots
are needed.  The bound on $t$ follows already with
$D(120)$ replaced by $.00000025$.
Lemma~\ref{per:coupled}, applied to~\eqref{per:seed-s}--\eqref{per:seed-t}
at $(n,b)=(60,b_1)$, now gives
\[
 1-S(120,b_2)<.052,\qquad 1-T(240,b_2)<.052.
\]

\emph{Induction.}
For $0\le u\le .052$,
\begin{equation}\label{per:G}
 f\bigl(1-(1-u)^3\bigr)=G(u)u^2,\qquad
 G(u)=\frac12\left(\sum_{r=1}^3(1-u)^{-r}\right)^2.
\end{equation}
The function $G$ is increasing, and exact substitution gives
$G(.052)<5.583<6$.  Also $6u^2\le u$ on this interval.
Thus, if the two defects in~\eqref{per:tracked-purities} are at most
$u_j$, the first coupled output is at most $6u_j^2=u_{j+1}\le u_j$.
Using this upper bound in the second output bounds that output by the
same number.  This proves~\eqref{per:tracked-purities} for every
$j\ge2$.  Each spatial purity exceeds $1/2$, so the maximal spatial
modulus remains unique and simple.  Applied to its probability list,
Lemma~\ref{per:purity} also gives
\[
 \sum_{i\ne0}\left|\frac{\mu_i(b_j)}{\lambda_j}\right|^{n_j}
 \le d(u_j),
\]
which proves~\eqref{per:spectral-ratio}.
The rational comparison $d(.052)<.971^{120}$ gives $\rho_2<.971$.

\emph{The sign and the rotation symmetry.}
At every level we have proved uniqueness of the eigenvalue of maximal
modulus.  If it were $-\lambda_j$, then
$\Tr(X_{b_j}^n)/\lambda_j^n$ would tend to $-1$ through odd integers
$n$.  Indeed, the residual ratios are bounded by some number below one,
and their fourth powers have finite sum, so their $n$th-power sum
vanishes as $n\to\infty$.  This contradicts the strictly positive
physical partition function $Z_n(b_j)$ at odd $n\ge5$.
The leading eigenvalue is therefore $+\lambda_j$.

Its eigenspace has a real unit vector, since $X_{b_j}$ is real.
The real unitaries $U_g$ commute with $X_{b_j}$, so they act on that
line by signs.
The order-three cyclic rotation $C$ must act by $+1$.  It conjugates
the three nonidentity elements of $D_2$ cyclically, so they have the
same sign on the line.  Since the product of any two is the third,
that common sign equals its square and hence is $+1$.
Thus every element of $\mathcal G=\langle D_2,C\rangle$ fixes the
leading eigenspace pointwise.

\emph{Comparing the leading eigenvalues.}
For every even $N\ge4$, physical squaring gives
$Z_N(b_{j+1})\le Z_N(b_j)^2$. The residual-tail estimate used above
gives $Z_N(b_j)/\lambda_j^N=1+o(1)$, so
$Z_N(b_j)^{1/N}\to\lambda_j$. Taking $N$th roots and letting $N$
tend to infinity through even integers therefore yields
$\lambda_{j+1}\le\lambda_j^2$.
For the opposite direction at the reference scale, recall that
$M_j=n_{j+1}$.  The leading spatial probability at $(M_j,b_{j+1})$
is at least $1-u_{j+1}$, so
\[
 \begin{split}
 \lambda_{j+1}^{M_j}
 &\ge(1-u_{j+1})Z_{M_j}(b_{j+1})\\
 &= (1-u_{j+1})T(M_j,b_j)Z_{M_j}(b_j)^2\\
 &\ge(1-u_{j+1})(1-u_j)\lambda_j^{2M_j}.
 \end{split}
\]
This proves~\eqref{per:lambda-ratios} and completes the proof.
\end{proof}
